\label{chap:related}
This chapter reviews the related work in the field of \gls{CRC} detection using \gls{DL} models. It covers the challenges faced by researchers and the various approaches taken to address these difficulties. Furthermore, it presents previous research that used the same \glspl{DS} used for training and validation in this work, as exposed in Chapters \ref{chap:methodology} and \ref{chap:results}. Finally, this chapter highlights the gaps in the existing literature and outlines how this dissertation aims to fill those gaps.

\section{Previous Colorectal Cancer Architectures}
\label{sec_prev_src_arch}

Before \glspl{CNN} became the dominant approach in the \gls{ML} field, handcrafted feature extractor techniques were the standard pattern for image recognition tasks. In 2012, the study by Xu et al. \cite{multipleClusterColon} approached the \gls{BC} of colon histopathology images via weakly \gls{SL}, while also identifying relevant tissue regions. Despite the absence of an \gls{ANN} approach, this research is considered one of the earliest \gls{ML} architectures dedicated to automated \gls{CRC} detection.

After the dissemination of the AlexNet architecture by Krizhevsky et al.~\cite{DBLP:conf/nips/KrizhevskySH12}, multiple studies were made in \gls{CNN} application fields. Ribeiro et al.~\cite{RibeiroDeepLearning} explored in 2016 the use of \gls{DL} models for robust feature extraction and efficient colonic polyp classification. The main difficulty of previous feature extraction methods in \gls{CRC} was the proper characterization of these patterns, which could vary due to multiple factors such as a lack or excess of illumination, blurring due to movement or water injection, and the appearance of the polyps themselves. Even with a limited quantity of image samples in the \gls{DS}, the pre-trained \gls{DL} model outperformed many previously diffused cancer classification systems and concluded that \glspl{CNN} were a good option for colonic polyp classification.

Korbar et al.~\cite{korbar2017deeplearningclassificationcolorectalpolyps} demonstrated the efficacy of ultra-deep \glspl{RNN} for the multi-class classification of colorectal polyps on whole-slide images. Their approach addressed the challenges of processing large-scale histological data by breaking the slides into smaller, overlapping patches. The study concludes that such automated tools can significantly reduce the cognitive burden on pathologists and improve the efficiency of colorectal cancer screening and risk assessment, achieving an accuracy of 93\%, precision of 89.7\%, and an F-score of 88.8\%.

\section{Previous Research on Selected Datasets}
\label{sec_dataset_usage}

Hossain \& Shaban \cite{DBLP:conf/aibthings/HossainS24} compared the results of \glspl{RNN} and the \gls{VGG}16 architecture using the \gls{LC25000} \gls{DS}. Although the primary objective of their article was to find a resource-efficient \gls{DL} architecture for \gls{CRC} detection, the results were extremely high, exceeding 99\% across all evaluated metrics \cite{metrics}. Given these exceptionally high scores, there is a valid concern regarding whether models trained exclusively on this \gls{DS} can generalize effectively to image sets from alternative sources, a problem that is explored in Chapter~\ref{chap:methodology}.

Kather et al.~\cite{kather2018crc100k} made available a \gls{DS} containing one hundred thousand histological images of human \gls{CRC} and healthy tissue. Researchers Tripathi et al.~\cite{colonML} utilized traditional \gls{ML} algorithms, such as \gls{KNN} and \gls{RFR}, to classify colon tissue, but their study lacks the inclusion of \gls{DL} models. In Chapters \ref{chap:methodology} and \ref{chap:results}, we address this gap by exploring \gls{DL} architectures trained on different available \glspl{DS} and validated with images from external sources~\cite{borkowski2019lungcoloncancerhistopathological, kather_2016_53169}.

Mulyadi \& Utaminingrum~\cite{10913097} explored different configurations of \glspl{RNN} combined with a self-attention layer for colon tissue classification using the \gls{CRC5000} \gls{DS} alongside ImageNet standard normalization \cite{DBLP:conf/nips/KrizhevskySH12}. However, their study did not include classic \gls{CNN} architectures for comparison.

\section{Quanvolution}
\label{sec_quanvolution_related_works}
Sobrinho et al.~\cite{quanvolution} explored the potential of \gls{QML} for breast cancer detection using the BreastMNIST \gls{DS}. Furthermore, the quantum circuits were validated with different image resolutions, varying from $28 \times 28$ pixels to $224 \times 224$ pixels during the training stage. The results suggest that quantum models may offer better generalization capabilities in handling resolution variations, with higher image resolutions not leading to a decline in classification performance.

The research also demonstrated an increase in training time for higher-resolution images and the capacity for generalization in breast cancer detection, despite the relatively small \gls{DS} used for training.

The conclusion of their work demonstrated that quantum models can achieve competitive performance compared to more complex, traditional \gls{CNN} models. This outcome raises important questions: Can this quantum solution be successfully expanded to other medical use cases, such as \gls{CRC} detection? Furthermore, can it be effectively applied to colon tissue detection?

\begin{figure}[htbp]
    \centering
    \includegraphics[width=0.9\textwidth]{mama-cancer-quanvolution}
    \caption{Illustration of the used hybrid quantum-classical model. Taken from \cite{quanvolution}.}
    \label{fig:breast_cancer_quanvolution}
\end{figure}


\section{Gaps Explored in This Work}
\label{sec_gaps_in_current_literature}
Most current research evaluates the performance of a trained model by splitting the \gls{DS}. However, these studies lack validation of generalization performance using alternative source images. Furthermore, most \gls{DL} studies present training metrics with only a single optimizer instead of comparing performance across multiple configurations.

Moreover, \gls{QML} is an emerging area that lacks contributions regarding colorectal histopathology image classification. Aside from the work of Sobrinho et al.~\cite{quanvolution} in breast cancer detection, \gls{QNN} applicability in colorectal polyp detection remains completely unexplored. Chapters~\ref{chap:methodology} and \ref{chap:results} exploit those gaps by using the same quantum circuit design as Sobrinho et al.~\cite{quanvolution} with the selected \glspl{DS} presented in this chapter.
